% =============================================================================
% Appendix A — VGC Bitmap & Sprite Shape RAM
% =============================================================================
\chapter{VGC Bitmap \& Sprite Shape RAM}

\epigraph{``The screen is a window into the machine's imagination.''}%
{\textit{--- Anonymous demo coder, 1989}}

\section*{Introduction}

The VGC manages two large memory spaces that live entirely outside the 6502's
64\,KB address space.  The graphics bitmap holds 64,000 bytes of pixel data for
the 320$\times$200 display, and the sprite shape RAM holds up to 256 shape
definitions of 256 bytes each (64\,KB total).  Neither space can be addressed
directly by the CPU; instead, the VGC provides two access mechanisms:

\begin{itemize}
  \item \textbf{MemRead/MemWrite commands} --- write a space selector and offset
        to the VGC command registers, then issue command \$19 (MemRead) or \$1A
        (MemWrite) at \addr{A010}.
  \item \textbf{DMA transfers} --- the DMA controller can copy blocks to and
        from VGC memory using space identifiers 3 (graphics bitmap) and 4
        (sprite shapes).
\end{itemize}


% ═══════════════════════════════════════════════════════════════════════════════
\section{Graphics Bitmap}
\label{sec:gfx-bitmap}
% ═══════════════════════════════════════════════════════════════════════════════

The graphics bitmap is a linear framebuffer of 320$\times$200 pixels stored in
row-major order.  Each pixel occupies one byte, with the colour index in the low
nibble (bits 3--0).  The upper nibble is ignored on write and reads as zero.

\begin{center}
\small
\begin{tabular}{@{}ll@{}}
\toprule
\textbf{Property}         & \textbf{Value} \\
\midrule
Resolution                 & 320$\times$200 pixels \\
Bits per pixel             & 4 (stored in low nibble of each byte) \\
Total size                 & 64,000 bytes \\
VGC MemRead/MemWrite space & \$02 (\texttt{MemSpaceGfx}) \\
DMA space                  & 3 (\texttt{DmaSpaceVgcGfx}) \\
\bottomrule
\end{tabular}
\end{center}

\subsection{Address Calculation}

The offset of any pixel within the bitmap is:

\[
  \textit{offset} = y \times 320 + x
\]

where $x \in [0, 319]$ and $y \in [0, 199]$.

\subsection{Row Start Offsets}

The following table shows the starting offset (in decimal and hex) for every
10th row:

\begin{center}
\small
\begin{tabular}{@{}rrr@{\qquad}rrr@{}}
\toprule
\textbf{Row} & \textbf{Offset (dec)} & \textbf{Offset (hex)} &
\textbf{Row} & \textbf{Offset (dec)} & \textbf{Offset (hex)} \\
\midrule
  0 &     0 & \$0000 &  100 & 32000 & \$7D00 \\
 10 &  3200 & \$0C80 &  110 & 35200 & \$8980 \\
 20 &  6400 & \$1900 &  120 & 38400 & \$9600 \\
 30 &  9600 & \$2580 &  130 & 41600 & \$A280 \\
 40 & 12800 & \$3200 &  140 & 44800 & \$AF00 \\
 50 & 16000 & \$3E80 &  150 & 48000 & \$BB80 \\
 60 & 19200 & \$4B00 &  160 & 51200 & \$C800 \\
 70 & 22400 & \$5780 &  170 & 54400 & \$D480 \\
 80 & 25600 & \$6400 &  180 & 57600 & \$E100 \\
 90 & 28800 & \$7080 &  190 & 60800 & \$ED80 \\
\bottomrule
\end{tabular}
\end{center}

\subsection{Drawing Commands vs.\ Direct Access}

Most programs use the VGC's built-in drawing commands (PLOT, LINE, CIRCLE, RECT,
FILL, PAINT) which operate on the bitmap internally.  The GCLS command clears
the entire bitmap to colour~0.  Direct MemRead/MemWrite access is useful for
reading individual pixel values or implementing custom rendering algorithms.

\subsection{BASIC Example: Read/Write Pixels via MemRead/MemWrite}

\begin{lstlisting}
10 REM read pixel at (100, 50)
20 POKE $A011, 0  : POKE $A012, 2  : REM space=gfx
30 OFFSET = 50*320+100
40 POKE $A013, OFFSET AND 255
50 POKE $A014, INT(OFFSET/256)
60 POKE $A010, $19 : REM MemRead
70 PRINT "COLOR:"; PEEK($A015)
80 REM
90 REM write pixel at (100, 50) to color 1 (white)
100 POKE $A015, 1
110 POKE $A010, $1A : REM MemWrite
\end{lstlisting}

\begin{notebox}
The VGC command registers at \addr{A011}--\addr{A014} retain their values after
a command executes.  When reading or writing consecutive pixels, you only need to
update the offset registers between commands.
\end{notebox}


% ═══════════════════════════════════════════════════════════════════════════════
\section{Sprite Shape RAM}
\label{sec:sprite-shapes}
% ═══════════════════════════════════════════════════════════════════════════════

Sprite shape RAM holds up to 256 shape definitions.  Each shape is a
16$\times$16 pixel image stored as 256 bytes (one byte per pixel, colour index
in the low nibble).  A pixel value of 0 is transparent by default.

\begin{center}
\small
\begin{tabular}{@{}ll@{}}
\toprule
\textbf{Property}         & \textbf{Value} \\
\midrule
Shape slots                & 256 (numbered 0--255) \\
Shape dimensions           & 16$\times$16 pixels \\
Bytes per shape            & 256 \\
Bits per pixel             & 4 (stored in low nibble of each byte) \\
Total size                 & 65,536 bytes (64\,KB) \\
VGC MemRead/MemWrite space & \$03 (\texttt{MemSpaceSprite}) \\
DMA space                  & 4 (\texttt{DmaSpaceVgcSprite}) \\
\bottomrule
\end{tabular}
\end{center}

\subsection{Address Calculation}

The offset of any pixel within sprite shape RAM is:

\[
  \textit{offset} = \textit{slot} \times 256 + \textit{row} \times 16 + \textit{column}
\]

where \textit{slot}~$\in [0, 255]$, \textit{row}~$\in [0, 15]$, and
\textit{column}~$\in [0, 15]$.

\subsection{Shape Slot Assignments}

\begin{notebox}
Shape slots 0--15 correspond to the default shapes for sprites 0--15 at
power-on, but any sprite can reference any shape slot by writing to the
\texttt{SprRegShape} register at offset~+4 within the sprite's register block
(base \addr{A040}).  This means multiple sprites can share a single shape
definition, and a sprite's appearance can be changed at any time by pointing it
to a different slot.
\end{notebox}

\subsection{BASIC Example: Define a Shape via MemWrite}

\begin{lstlisting}
10 REM define a simple arrow shape in slot 0
20 POKE $A011, 0 : POKE $A012, 3 : REM space=sprite
30 FOR ROW = 0 TO 15
40   FOR COL = 0 TO 15
50     OFFSET = 0*256 + ROW*16 + COL
60     POKE $A013, OFFSET AND 255
70     POKE $A014, INT(OFFSET/256)
80     C = 0 : REM transparent
90     IF COL = 7 AND ROW < 12 THEN C = 1
100    IF ROW = 0 AND COL >= 5 AND COL <= 9 THEN C = 1
110    POKE $A015, C
120    POKE $A010, $1A : REM MemWrite
130  NEXT COL
140 NEXT ROW
\end{lstlisting}

\subsection{Using DMA for Bulk Shape Loading}

For loading many shapes at once (e.g., from XRAM or CPU RAM), the DMA controller
is far more efficient than individual MemWrite commands.  Set the DMA source
space to CPU RAM (space~0) or XRAM (space~5), the destination space to
\texttt{DmaSpaceVgcSprite} (space~4), and issue a DMA copy command.  See
Chapter~\ref{sec:xmc-registers} for DMA register details.

\begin{lstlisting}
10 REM DMA copy 2560 bytes (10 shapes) from CPU RAM
20 REM at $2000 into sprite shape slot 0
30 DMACOPY $2000, 0, 2560, 0, 4
\end{lstlisting}
